\documentclass[11pt,letterpaper]{article}

\usepackage[margin=1in]{geometry}
\usepackage{amsmath,amssymb}
\usepackage{booktabs}
\usepackage{array}
\usepackage{graphicx}
\usepackage{enumitem}
\usepackage{tabularx}
\usepackage{titlesec}
\usepackage[hidelinks]{hyperref}

\titleformat{\section}{\Large\bfseries}{\thesection}{1em}{}
\titleformat{\subsection}{\large\bfseries}{\thesubsection}{1em}{}
\setlength{\parskip}{0.5em}
\setlength{\parindent}{0pt}

\title{\textbf{ECE332 Lab 3: Waveguides}}
\author{Bao Nguyen, Nick Stites}
\date{\textbf{Term:} Spring 2026}

\begin{document}

\maketitle
\thispagestyle{empty}
\newpage

\section*{Pre-Lab Assignment}

\textbf{Reading:} Sections 8-6 through 8-11 in Ulaby.\\
\textbf{Equipment:} Systron/Donner 5000 Sweeper/Oscillator and HP 415E SWR Meter.

\subsection*{Answers}

\textbf{Question 1.} Which transmission lines do NOT support TEM waves? (all that apply)\\
\textit{Answer:} (a) Rectangular Waveguides and (d) Optical Fiber.

\textbf{Question 2.} Dominant mode for transverse magnetic (TM) waves?\\
\textit{Answer:} (d) TM$_{11}$.

\textbf{Question 3.} Dominant mode for transverse electric (TE) waves?\\
\textit{Answer:} (b) TE$_{10}$.

\textbf{Question 4.} What type of filter does a resonant cavity implement?\\
\textit{Answer:} (d) Bandpass.

\textbf{Question 5.} Which instrument measures the Standing Wave Ratio (SWR)?\\
\textit{Answer:} (b) Hewlett Packard 415E.

\newpage

\section*{Experiment 1: Wavelength Measurements}

\subsection*{Objective}

Measure the standing-wave pattern inside a slotted rectangular waveguide at three input frequencies (11~GHz, 9.5~GHz, 8.5~GHz) using a traveling probe and SWR meter. From the minima/maxima positions calculate the waveguide wavelength $\lambda_g$, phase velocity $v_p$, and group velocity $v_g$ for each frequency. Also measure the physical dimensions $a$ and $b$ of the waveguide and use them to compute the cutoff frequencies of the three lowest modes.

\subsection*{Theory}

In a two-conductor transmission line, the electric field is transverse to the direction of propagation, giving rise to TEM waves. A hollow metallic waveguide has only one conductor, so TEM waves cannot propagate. Instead, TE and TM modes exist, each with a minimum (cutoff) frequency below which the mode is evanescent.

For a rectangular waveguide with inner width $a$ (along $x$) and height $b$ (along $y$), the cutoff frequency of the TE$_{mn}$ / TM$_{mn}$ mode is:
\begin{equation}
f_c = \frac{1}{2\sqrt{\mu\varepsilon}}\sqrt{\left(\frac{m}{a}\right)^2 + \left(\frac{n}{b}\right)^2}
\end{equation}

The dominant mode is TE$_{10}$ ($m=1$, $n=0$) with $f_c = c/(2a)$. For propagation to occur, the operating frequency $f$ must exceed $f_c$.

Inside the waveguide the wave bounces at an angle between the walls, producing a standing wave along $z$ when the far end is short-circuited. The waveguide wavelength $\lambda_g$ is longer than the free-space wavelength $\lambda = c/f$:
\begin{equation}
\lambda_g = \frac{2\pi}{\beta} = \frac{2\pi}{\sqrt{\omega^2\mu\varepsilon - (m\pi/a)^2 - (n\pi/b)^2}}
\end{equation}

The phase velocity (speed of the phase front along $z$) is:
\begin{equation}
v_p = f\lambda_g = \frac{v_{p0}}{\sqrt{1 - (f_c/f)^2}}
\end{equation}
Note $v_p > c$; it carries no information. The group velocity (information/energy speed) is:
\begin{equation}
v_g = v_{p0}\sqrt{1 - (f_c/f)^2} = \frac{1}{\sqrt{\mu\varepsilon}}\cdot\frac{\lambda}{\lambda_g}
\end{equation}
Note $v_p \cdot v_g = c^2$ (for air-filled guide), so $v_g < c$.

The waveguide wavelength is measured from successive minima/maxima. Each consecutive min--max pair is $\lambda_g/4$ apart:
\begin{equation}
\lambda_g = 4\,\Delta d = 4(d_{n+1} - d_n)
\end{equation}

\subsection*{Setup}

Equipment used:
\begin{itemize}[leftmargin=*]
  \item Systron/Donner 5000 Sweeper/Oscillator with waveguide cables (8--12.4~GHz)
  \item Slotted waveguide section with traveling probe on carriage (vernier scale)
  \item Short-circuit termination
  \item SWR Meter (HP 415E) connected via BNC cable
  \item Meter stick (for physical dimension measurement)
\end{itemize}

The short-circuit termination was attached to the output of the slotted waveguide. The probe was connected to the SWR meter. The sweeper/oscillator was set to 11~GHz (2F-C knob), RF Level 5--10~mW, ``1~kHz'' switch ON. Probe penetration depth was tuned for mid-scale deflection (RANGE-DB 30--40~dB). Measurements were repeated for 9.5~GHz and 8.5~GHz.

\subsection*{Equations Used}

\begin{itemize}[leftmargin=*]
  \item $f_c = \dfrac{1}{2\sqrt{\mu\varepsilon}}\sqrt{(m/a)^2 + (n/b)^2}$ -- cutoff frequency (Eq.~1)
  \item $\lambda_g = 2\pi/\beta$ -- waveguide wavelength (Eq.~2)
  \item $v_p = f\lambda_g = v_{p0}/\sqrt{1-(f_c/f)^2}$ -- phase velocity (Eq.~3)
  \item $v_g = v_{p0}\sqrt{1-(f_c/f)^2} = c^2/v_p$ -- group velocity (Eq.~4)
  \item $\lambda_g = 4\Delta d$ from consecutive min/max positions (Eq.~5)
\end{itemize}

\subsection*{Results / Data}

\textbf{Table 1.} Waveguide SWR minima/maxima positions and calculated quantities. Note: $\lambda_g$, $v_p$, $v_g$ are calculated only for Min rows using $\lambda_g = 4\times|\text{Min}_n - \text{Max}_n|$. The Average row uses all 5 consecutive $\Delta d$ pairs.

\begin{center}
\begin{tabular}{lcccc}
\toprule
 & Position (mm) & $\lambda_g$ (mm) & $v_p$ (m/s) & $v_g$ (m/s) \\
\midrule
\multicolumn{5}{l}{\textit{$f = 11$~GHz}} \\
Min 1     & 170.1 & 31.2 & $3.432\times10^{8}$ & $2.622\times10^{8}$ \\
Max 1     & 162.3 & 36.0 & $3.960\times10^{8}$ & $2.273\times10^{8}$ \\
Min 2     & 153.3 & 32.4 & $3.564\times10^{8}$ & $2.525\times10^{8}$ \\
Max 2     & 145.2 & 34.4 & $3.784\times10^{8}$ & $2.378\times10^{8}$ \\
Min 3     & 136.6 & 29.6 & $3.256\times10^{8}$ & $2.764\times10^{8}$ \\
Max 3     & 129.2 &      &                     &                     \\
Average:  &       & 32.7 & $3.599\times10^{8}$ & $2.501\times10^{8}$ \\
\midrule
\multicolumn{5}{l}{\textit{$f = 9.5$~GHz}} \\
Min 1     & 174.4 & 44.0 & $4.180\times10^{8}$ & $2.153\times10^{8}$ \\
Max 1     & 163.4 & 43.6 & $4.142\times10^{8}$ & $2.173\times10^{8}$ \\
Min 2     & 152.5 & 42.0 & $3.990\times10^{8}$ & $2.256\times10^{8}$ \\
Max 2     & 142.0 & 42.4 & $4.028\times10^{8}$ & $2.234\times10^{8}$ \\
Min 3     & 131.4 & 43.6 & $4.142\times10^{8}$ & $2.173\times10^{8}$ \\
Max 3     & 120.5 &      &                     &                     \\
Average:  &       & 43.1 & $4.096\times10^{8}$ & $2.197\times10^{8}$ \\
\midrule
\multicolumn{5}{l}{\textit{$f = 8.5$~GHz}} \\
Min 1     & 166.4 & 62.4 & $5.304\times10^{8}$ & $1.697\times10^{8}$ \\
Max 1     & 150.8 & 47.6 & $4.046\times10^{8}$ & $2.224\times10^{8}$ \\
Min 2     & 138.9 & 64.4 & $5.474\times10^{8}$ & $1.644\times10^{8}$ \\
Max 2     & 122.8 & 46.0 & $3.910\times10^{8}$ & $2.302\times10^{8}$ \\
Min 3     & 111.3 & 58.0 & $4.930\times10^{8}$ & $1.826\times10^{8}$ \\
Max 3     & 96.8  &      &                     &                     \\
Average:  &       & 55.7 & $4.733\times10^{8}$ & $1.902\times10^{8}$ \\
\bottomrule
\end{tabular}
\end{center}

Measured waveguide dimensions:
\begin{itemize}[leftmargin=*]
  \item $a = 22.92$~mm
  \item $b = 10.29$~mm
\end{itemize}

Distance from short circuit to nearest detectable minimum:
\begin{itemize}[leftmargin=*]
  \item 11~GHz: 83.5~mm
  \item 9.5~GHz: 95.6~mm
  \item 8.5~GHz: 85.4~mm
\end{itemize}

\subsection*{Discussion}

\paragraph{Q1. Explain the difference between $\lambda$ and $\lambda_g$.}
The free-space wavelength $\lambda = c/f$ is the spatial period of the wave if it propagated in a straight line through free space at the speed of light. Inside the waveguide, the TE$_{10}$ mode can be visualized as two plane waves bouncing obliquely between the broad walls. These waves travel at an angle $\theta$ to the $z$-axis, so their axial phase advances more slowly than in free space. The resulting axial period, namely the waveguide wavelength $\lambda_g = \lambda/\sqrt{1-(f_c/f)^2}$, is always greater than $\lambda$ and grows to infinity as $f \to f_c$.

\paragraph{Q2. Electric field at the short-circuit termination.}
A short-circuit termination enforces the boundary condition that the tangential electric field is zero at the conducting end wall: $E_\text{tan} = 0$. The probe measures the tangential (transverse) E-field. Therefore, if the probe could be positioned exactly at the short circuit it would read zero signal. The short is a node of the standing wave.

\paragraph{Q3. Expected distance from short circuit to nearest detectable minimum.}
The short circuit forces $E = 0$ at $z = 0$. Minima repeat every $\lambda_g/2$. The nearest detectable minimum is $n\times(\lambda_g/2)$ from the termination, where $n$ is the smallest integer that places the minimum within probe reach.

\begin{center}
\begin{tabular}{cccccc}
\toprule
$f$ (GHz) & $\lambda_g/2$ (mm) & $n$ & Expected $d$ (mm) & Measured $d$ (mm) & \% diff \\
\midrule
11   & 16.36 & 5            & 81.80 & 83.5 & 2.1\%  \\
9.5  & 21.56 & 4 ($\approx4.43$) & 86.24 & 95.6 & 10.9\% \\
8.5  & 27.84 & 3            & 83.52 & 85.4 & 2.3\%  \\
\bottomrule
\end{tabular}
\end{center}

The 11~GHz and 8.5~GHz results agree within about 2\%, consistent with typical vernier scale reading error. The 9.5~GHz result shows a larger discrepancy (about 10\%), likely due to probe alignment or a measurement made to a different minimum than expected.

\paragraph{Q4. Half-wavelengths between nearest detectable minimum and short circuit.}
$n = d_\text{measured} / (\lambda_g/2)$:

\begin{center}
\begin{tabular}{ccccc}
\toprule
$f$ (GHz) & Measured $d$ (mm) & $\lambda_g/2$ (mm) & $n = d/(\lambda_g/2)$ & Nearest int $n$ \\
\midrule
11   & 83.5 & 16.36 & 5.10 & 5 \\
9.5  & 95.6 & 21.56 & 4.43 & 4 (with uncertainty) \\
8.5  & 85.4 & 27.84 & 3.07 & 3 \\
\bottomrule
\end{tabular}
\end{center}

\paragraph{Q5. Cutoff frequencies and propagating modes.}
(a) Using $a = 22.92$~mm, $b = 10.29$~mm, $c = 3\times10^{8}$~m/s (Eq.~1):

Only TE modes are relevant for the three lowest. TM requires $m \geq 1$ AND $n \geq 1$, so TM$_{11}$ is always higher than the lowest TE modes.

\begin{center}
\begin{tabular}{cccl}
\toprule
Mode & $(m,n)$ & $f_c$ formula & $f_c$ (GHz) \\
\midrule
TE$_{10}$         & $(1,0)$ & $c/(2a)$                  & 6.54 (lowest)        \\
TE$_{20}$         & $(2,0)$ & $c/a$                     & 13.09                \\
TE$_{01}$         & $(0,1)$ & $c/(2b)$                  & 14.58                \\
TE$_{11}$/TM$_{11}$ & $(1,1)$ & $(c/2)\sqrt{1/a^2+1/b^2}$ & 15.98 (4th, not in top 3) \\
\bottomrule
\end{tabular}
\end{center}

Three lowest: TE$_{10}$ (6.54~GHz) $<$ TE$_{20}$ (13.09~GHz) $<$ TE$_{01}$ (14.58~GHz). Note $a/b = 22.92/10.29 = 2.23 \approx 2$, confirming standard X-band ordering TE$_{10} <$ TE$_{20} <$ TE$_{01}$.

(b) Propagating modes ($f > f_c$ required):

All three operating frequencies (8.5, 9.5, 11~GHz) lie between TE$_{10}$ cutoff (6.54~GHz) and TE$_{20}$ cutoff (13.09~GHz). Therefore only TE$_{10}$ propagates at all three frequencies, so the guide operates single-mode throughout.

(c) Back-calculating $f_c$ from measured $\lambda_g$ using $f_c = f\sqrt{1 - (\lambda/\lambda_g)^2}$:

\begin{center}
\begin{tabular}{cccc}
\toprule
Frequency (GHz) & $\lambda$ (mm) & $\lambda_g$ (mm) & $f_c$ (GHz) \\
\midrule
11   & 27.27 & 32.73 & 6.08 \\
9.5  & 31.58 & 43.12 & 6.47 \\
8.5  & 35.29 & 55.68 & 6.57 \\
\bottomrule
\end{tabular}
\end{center}

Average back-calculated $f_c = 6.37$~GHz vs geometric $f_c$ (TE$_{10}$) = 6.54~GHz, so the percent difference is 2.6\%. Agreement is within about 3\%, consistent with minor probe loading and vernier scale reading error.

\newpage

\section*{Experiment 2: Frequency Measurements}

\subsection*{Objective}

Use a tunable resonant cavity (HP X532B frequency meter) to independently measure the microwave frequency at each of the three oscillator settings (11~GHz, 9.5~GHz, 8.5~GHz). Compare the cavity-measured frequency to the oscillator dial setting and recalculate phase and group velocities.

\subsection*{Theory}

A resonant cavity is a hollow metallic enclosure in which EM fields oscillate at discrete resonant frequencies determined by its dimensions. When a signal at the resonant frequency passes through the cavity it is strongly absorbed, causing a pronounced dip in the SWR meter reading. Tuning the cavity until the SWR dip is greatest gives the signal frequency directly from the calibrated scale.

Combining the cavity-measured frequency with $\lambda_g$ from Experiment 1 allows an independent cross-check of phase velocity:
\begin{align}
v_p &= \lambda_g \cdot f \\
v_g &= c^2 / v_p
\end{align}

\subsection*{Setup}

Equipment: same as Experiment 1 with the addition of HP X532B Tunable Cavity Frequency Meter inserted between the carriage and the short-circuit termination. The probe was positioned at a standing-wave maximum. The cavity knob was tuned for maximum SWR dip at each oscillator frequency setting.

\subsection*{Equations Used}

\begin{itemize}[leftmargin=*]
  \item $v_p = \lambda_g\cdot f$ -- phase velocity
  \item $v_g = c^2 / v_p$ -- group velocity
  \item \% difference $= |\text{measured} - \text{expected}| / \text{expected} \times 100\%$
\end{itemize}

\subsection*{Results / Data}

\textbf{Table 2.} Resonant cavity frequency measurements and derived velocities.

\begin{center}
\begin{tabular}{ccccc}
\toprule
Oscillator $f$ (GHz) & Measured $f$ (GHz) & $\lambda_g$ avg (mm) & $v_p$ (m/s) & $v_g$ (m/s) \\
\midrule
11   & 11.050 & 32.7 & $3.616\times10^{8}$ & $2.489\times10^{8}$ \\
9.5  & 9.405  & 43.1 & $4.055\times10^{8}$ & $2.219\times10^{8}$ \\
8.5  & 8.515  & 55.7 & $4.741\times10^{8}$ & $1.898\times10^{8}$ \\
\bottomrule
\end{tabular}
\end{center}

\subsection*{Discussion}

\paragraph{Q1. Phase velocity: $v_p = \lambda_g \cdot f$.}
Using the measured values for frequency and wavelength, calculate the phase velocity of the wave: $v_p = \lambda_g\cdot f$. How close (percent difference) is this to the values from Experiment~1? Based on this, how important is it to know the frequency accurately?

Using cavity-measured frequency and average $\lambda_g$ from Experiment 1:
\begin{itemize}[leftmargin=*]
  \item $v_p(11~\text{GHz}) = 32.72~\text{mm} \times 11.050~\text{GHz} = 3.616\times10^{8}$~m/s (Exp 1: $3.599\times10^{8}$~m/s, \% diff $= 0.5\%$)
  \item $v_p(9.5~\text{GHz}) = 43.12~\text{mm} \times 9.405~\text{GHz} = 4.055\times10^{8}$~m/s (Exp 1: $4.096\times10^{8}$~m/s, \% diff $= 1.0\%$)
  \item $v_p(8.5~\text{GHz}) = 55.68~\text{mm} \times 8.515~\text{GHz} = 4.741\times10^{8}$~m/s (Exp 1: $4.733\times10^{8}$~m/s, \% diff $= 0.2\%$)
\end{itemize}

All three values agree with Experiment 1 to within 1.1\%, confirming that the oscillator frequency was close to its set value and that the waveguide wavelength measurements were accurate. Small discrepancies are attributable to vernier scale reading error and minor oscillator frequency drift.

Based on this, knowing the frequency accurately is important but not critical at this level. Since $v_p = \lambda_g \cdot f$, a 1\% error in frequency directly produces a 1\% error in $v_p$. More critically, frequency accuracy matters for mode control: if the operating frequency drifted below the TE$_{20}$ cutoff (13.09~GHz) it could excite unwanted higher-order modes and corrupt measurements. The cavity meter confirms the oscillator stays within about 0.5\% of its set value, which is sufficient for single-mode waveguide operation.

\paragraph{Q2. What happens to $v_g$ as $f$ approaches $f_c$?}
From Eq.~4: $v_g = v_{p0}\sqrt{1 - (f_c/f)^2}$. As $f \to f_c$, the term $(f_c/f)^2 \to 1$, so $\sqrt{1 - (f_c/f)^2} \to 0$ and therefore $v_g \to 0$. The measured data confirms this trend:

\begin{center}
\begin{tabular}{cccc}
\toprule
$f$ (GHz) & $f_c$ (GHz) & $f / f_c$ & $v_g$ (m/s) \\
\midrule
11   & 6.54 & 1.68 & $2.501\times10^{8}$ \\
9.5  & 6.54 & 1.45 & $2.197\times10^{8}$ \\
8.5  & 6.54 & 1.30 & $1.902\times10^{8}$ \\
\bottomrule
\end{tabular}
\end{center}

As frequency decreases toward $f_c = 6.54$~GHz, $v_g$ drops from $2.501\times10^{8}$ to $1.902\times10^{8}$~m/s, approaching zero at cutoff.

Physical explanation: the TE$_{10}$ mode can be visualized as two plane waves bouncing obliquely between the broad walls at angle $\theta$ to the $z$-axis, where $\sin\theta = f_c/f$. As $f \to f_c$, $\theta \to 90^\circ$, so the waves bounce nearly perpendicular to the axis with almost no net axial advance. Energy is trapped bouncing transversely and travels axially more and more slowly. At exactly $f = f_c$ the waves bounce straight across ($\theta = 90^\circ$) and $v_g = 0$, meaning the mode carries no energy along the guide. Below cutoff the mode becomes evanescent and decays exponentially.

\end{document}
